\documentclass[aspectratio=169]{beamer}

% -------------------------------------------------
% Tema e cores
% -------------------------------------------------
\usetheme{Madrid}
\usecolortheme{default}

\definecolor{mainblue}{HTML}{1F4E79}
\definecolor{darkgray}{HTML}{444444}

\setbeamercolor{title}{fg=mainblue}
\setbeamercolor{frametitle}{fg=mainblue}
\setbeamercolor{structure}{fg=mainblue}
\setbeamercolor{normal text}{fg=darkgray}

\setbeamertemplate{navigation symbols}{}

% -------------------------------------------------
% Pacotes
% -------------------------------------------------
\usepackage[brazil]{babel}
\usepackage[utf8]{inputenc}
\usepackage{graphicx}
\usepackage{booktabs}

% -------------------------------------------------
% Título
% -------------------------------------------------
\title[Revisão Sistemática]{
Metodologia de Investigação em Revisão Sistemática
}
\subtitle{
Tema 1 – Apontamento de Literatura
}
\author{
Disciplina: Revisão Sistemática e Meta-análises\\
Pós-graduação (Mestrado e Doutoramento)
}
\date{}

% -------------------------------------------------
\begin{document}

% -------------------------------------------------
\begin{frame}
\titlepage
\end{frame}

% =================================================
\begin{frame}{Introdução à Investigação Científica}
% =================================================

A investigação científica constitui um processo sistemático, rigoroso e cumulativo de produção de conhecimento, orientado pela observação empírica, formulação teórica e validação metodológica.

\vspace{0.3cm}

No contexto da revisão sistemática, a investigação assume um papel central na organização, síntese e avaliação crítica da evidência existente.

\end{frame}

% =================================================
\begin{frame}{Objetivos da Investigação Científica}
% =================================================

Os principais objetivos da investigação científica incluem:

\begin{itemize}
    \item Descrever fenómenos naturais ou sociais;
    \item Explicar relações causais ou associativas entre variáveis;
    \item Prever comportamentos ou resultados futuros;
    \item Sustentar a tomada de decisão baseada na evidência.
\end{itemize}

\vspace{0.3cm}

A pesquisa científica distingue-se de abordagens narrativas ou opinativas pela sua reprodutibilidade, transparência metodológica e fundamentação empírica.

\end{frame}

% =================================================
\begin{frame}{Pirâmide dos Níveis de Evidência}
% =================================================

A hierarquização dos níveis de evidência é fundamental para a avaliação da robustez científica dos resultados de investigação.

\begin{itemize}
    \item Diferentes desenhos produzem diferentes níveis de evidência;
    \item O controlo de viés aumenta à medida que se sobe na pirâmide;
    \item Revisões sistemáticas sintetizam estudos independentes.
\end{itemize}

\end{frame}

% =================================================
\begin{frame}{Pirâmide dos Níveis de Evidência}
% =================================================

\begin{center}
\includegraphics[width=0.6\textwidth]{Tema-1/fig1_quant_vs_qual.png}
\end{center}

\vspace{0.2cm}

Revisões sistemáticas e meta-análises ocupam o nível mais elevado da hierarquia da evidência científica.

\end{frame}

% =================================================
\begin{frame}{Tipos de Desenhos de Investigação}
% =================================================

Os desenhos de investigação podem ser classificados de acordo com a natureza dos dados e os objetivos analíticos:

\begin{itemize}
    \item Investigação quantitativa;
    \item Investigação qualitativa;
    \item Investigação de métodos mistos.
\end{itemize}

Estes desenhos determinam critérios de inclusão, estratégias de síntese e métodos de análise em revisões sistemáticas.

\end{frame}

% =================================================
\begin{frame}{Desenhos de Investigação}
% =================================================

\begin{table}
\centering
\begin{tabular}{p{3cm} p{5cm} p{4cm}}
\toprule
\textbf{Tipo} & \textbf{Características} & \textbf{Exemplos} \\
\midrule
Quantitativo & Dados numéricos, inferência estatística & Ensaios clínicos, coortes \\
Qualitativo & Interpretação de fenómenos complexos & Entrevistas, grupos focais \\
Misto & Integração de abordagens & Surveys + entrevistas \\
\bottomrule
\end{tabular}
\end{table}

\end{frame}

% =================================================
\begin{frame}{Elementos Fundamentais de um Estudo}
% =================================================

Um estudo científico robusto deve contemplar:

\begin{enumerate}
    \item Formulação clara do problema de pesquisa;
    \item Definição explícita de objetivos e hipóteses;
    \item Identificação das variáveis relevantes;
    \item Metodologia transparente e reprodutível;
    \item Estratégia de análise adequada;
    \item Interpretação crítica dos resultados.
\end{enumerate}

\end{frame}

% =================================================
\begin{frame}{Medidas de Força em Investigação Quantitativa}
% =================================================

As medidas de força quantificam associações, diferenças ou efeitos entre variáveis.

\vspace{0.3cm}

Principais medidas utilizadas em meta-análises:

\begin{itemize}
    \item Risco Relativo (RR);
    \item Odds Ratio (OR);
    \item Coeficiente de correlação;
    \item Tamanho do efeito (Cohen’s d, Hedges’ g).
\end{itemize}

\end{frame}

% =================================================
\begin{frame}{Investigação Quantitativa vs. Qualitativa}
% =================================================

\begin{table}
\centering
\begin{tabular}{p{4cm} p{4cm} p{4cm}}
\toprule
\textbf{Aspecto} & \textbf{Quantitativa} & \textbf{Qualitativa} \\
\midrule
Tipo de dados & Numéricos & Textuais ou visuais \\
Objetivo & Testar hipóteses & Explorar significados \\
Análise & Estatística & Interpretativa \\
\bottomrule
\end{tabular}
\end{table}

\end{frame}

% =================================================
\begin{frame}{Comparação Conceptual}
% =================================================

\begin{center}
\includegraphics[width=0.7\textwidth]{Tema-1/fig2_piramide_evidencia.png}
\end{center}

\end{frame}

% =================================================
\begin{frame}{Avaliação Crítica de um Artigo Científico}
% =================================================

A avaliação crítica constitui uma competência central em revisões sistemáticas.

\begin{enumerate}
    \item A questão de investigação está claramente definida?
    \item O desenho do estudo é apropriado?
    \item A amostra é adequada?
    \item Os métodos são válidos e confiáveis?
    \item A análise estatística é apropriada?
\end{enumerate}

\end{frame}

% =================================================
\begin{frame}{Avaliação Crítica (continuação)}
% =================================================

\begin{enumerate}
    \setcounter{enumi}{5}
    \item Os resultados são apresentados de forma clara?
    \item As conclusões são sustentadas pelos dados?
    \item As limitações são discutidas?
    \item O estudo contribui para o avanço do conhecimento?
\end{enumerate}

\vspace{0.3cm}

Esta estrutura orienta decisões de inclusão e exclusão em revisões sistemáticas.

\end{frame}

% -------------------------------------------------
\end{document}
